% Problemas del capítulo 7. Generado por tools/problemas/gen.py; no editar a mano.
\section{Problemas}\label{sec:problemas-07}

Graduados: \({\star}\) directo, \({\star}{\star}\) para pensar, \({\star}{\star}{\star}\) desafío. Las respuestas están en el apéndice~\ref{ap:respuestas}; conviene intentarlos antes de mirarlas.

\begin{problema}[\({\star}\)]\label{pb:07-1}
Encontrá los autovalores y autovectores de \(A=\begin{pmatrix}3&1\\0&2\end{pmatrix}\).
\end{problema}

\begin{problema}[\({\star}{\star}\)]\label{pb:07-2}
Hacé lo mismo con la simétrica \(S=\begin{pmatrix}5&2\\2&2\end{pmatrix}\) y comprobá que los autovectores son perpendiculares.
\end{problema}

\begin{problema}[\({\star}{\star}\)]\label{pb:07-3}
Mostrá que el giro de \(30^\circ\) del plano no tiene autovalores reales. ¿Qué significa eso geométricamente?
\end{problema}

\begin{problema}[\({\star}{\star}\)]\label{pb:07-4}
La matriz que lleva \(e_1\to e_2\), \(e_2\to e_3\) y \(e_3\to e_1\) es un giro del espacio. Encontrá su eje (el autovector de autovalor \(1\)) y el ángulo de giro.
\end{problema}

\begin{problema}[\({\star}{\star}{\star}\)]\label{pb:07-5}
En una placa se mide el tensor de deformación \(E=\begin{pmatrix}-200&150\\150&400\end{pmatrix}\) (en microdeformaciones, \(\mu\varepsilon\)). Calculá las deformaciones principales \(e_1\), \(e_2\) y el ángulo de la dirección principal con el eje \(x\).
\end{problema}
